\medterm{Kabuki Syndrome} A rare genetic disorder that may cause distinctive facial features, short stature, developmental or learning difficulties, weak muscle tone, skeletal changes, and heart defects. Hearing loss, seizures, immune problems, and kidney abnormalities may also occur. Most cases result from a change in the \textit{KMT2D} or \textit{KDM6A} gene.

\synonyms
Niikawa-Kuroki Syndrome

\medterm{Kahler Disease} Another name for multiple myeloma, a cancer of plasma cells in the bone marrow. It can produce an abnormal protein and cause weak or damaged bones, high calcium, kidney problems, and anemia.
\synonyms
Multiple Myeloma; Plasma Cell Myeloma

\medterm{Kahler's Disease} Alternate name; see \textit{Multiple Myeloma}.

\medterm{Kala-Azar} A serious parasitic infection caused by Leishmania species and transmitted by sandflies. It may cause prolonged fever, weight loss, enlargement of the spleen and liver, and suppression of blood-cell production.

\synonyms
Black Sickness; Dumdum Fever

\medterm{Kallikrein-Kinin System} A group of blood and tissue proteins that makes chemical messengers such as bradykinin. These messengers widen small blood vessels, make them more leaky, and contribute to pain and inflammation.
\synonyms
Kininergic System; Bradykinin Cascade

\medterm{Kallmann Syndrome} An inherited condition that causes delayed or absent puberty and a reduced or absent sense of smell. It occurs when the brain does not send the hormone signals needed to start puberty.

\medterm{Kallmann Syndrome and IHH} Conditions in which deficient hypothalamic or pituitary gonadotropin signaling causes low sex hormones and absent or delayed puberty. Kallmann syndrome also includes impaired smell and may be genetic; fertility and sex-hormone replacement require endocrine care.

\medterm{Kanamycin} An antibiotic used for some serious bacterial infections. It can permanently damage hearing and harm the kidneys, especially with prolonged or high-dose use.

\medterm{Kangaroo Care} Holding a newborn directly against a caregiver’s bare chest while monitoring the infant’s
condition. It supports thermal regulation, physiologic stability, bonding, and feeding,
and is particularly useful for many preterm or low-birth-weight infants.

\synonyms
Kangaroo Mother Care; KMC

\medterm{Kaolin} A soft clay mineral used in some medicines, laboratory products, and industrial materials. Breathing large amounts of kaolin dust over time can damage the lungs.

\medterm{Kaposiform Hemangioendothelioma} A rare blood-vessel tumor that usually affects infants and young children. It can grow into nearby tissue and may cause a serious loss of platelets and problems with blood clotting.

\synonyms
KHE

\medterm{Kaposi Sarcoma} A vascular neoplasm associated with human herpesvirus 8 (HHV-8). It may cause nonblanching red, purple, brown, or black macules, plaques, or nodules on the skin or oral mucosa and can involve lymph nodes or internal organs; disease is more likely or aggressive with immunosuppression, including HIV or transplantation.

\medterm{Kaposi's Sarcoma} A vascular tumor associated with human herpesvirus 8, also called Kaposi sarcoma-associated herpesvirus. It may cause purple, red, or brown lesions of the skin or oral mucosa and can involve lymph nodes or internal organs; it is more likely or aggressive when immunity is weakened, such as with HIV infection or after organ transplantation.

\medterm{Karnofsky Performance Status} A scale used to describe how well a person can perform daily activities and care for themselves. It helps clinicians assess overall function and ability to tolerate treatment, especially in cancer care.

\medterm{Kartagener Syndrome} An inherited condition in which tiny hair-like structures that clear mucus do not work properly. It causes repeated sinus and lung infections and may place the internal organs in a reversed position.

\medterm{Kartagener Triad} The combination of reversed internal-organ positions, long-term sinus inflammation, and widened damaged airways. It results from an inherited problem with microscopic moving hairs that clear mucus.

\medterm{Karyomegaly} Abnormal enlargement of a cell’s nucleus, the part that contains most of its genetic material. It may occur with injury, infection, toxic exposure, tissue repair, or cancer and is interpreted with other findings.

\medterm{Karyotype} An organized picture of a person’s chromosomes, arranged in pairs to show their number and large-scale structure. It can detect some missing, extra, or rearranged chromosomes but not every genetic condition.

\medterm{Karyotype Result} A laboratory report describing the number and visible structure of a person’s chromosomes. It may identify some chromosome differences, such as an extra chromosome, but a normal result does not exclude all genetic disorders.

\medterm{Karyotyping} A laboratory test that examines chromosomes in dividing cells to detect some missing, extra, or rearranged chromosomes. It may be used in pregnancy, infertility, repeated miscarriage, developmental conditions, or blood cancers.

\medterm{Kasabach-Merritt Phenomenon} A serious blood-clotting problem linked to certain fast-growing blood-vessel tumors in babies. Platelets become trapped in the tumor, causing a very low platelet count, anemia from broken red blood cells, and low levels of clotting proteins.

\synonyms
KMP; Kasabach-Merritt Syndrome

\medterm{Kasabach-Merritt Syndrome} Another name for Kasabach-Merritt phenomenon: a serious clotting problem in babies with certain fast-growing blood-vessel tumors. It causes very low platelets, anemia from broken red blood cells, and low levels of clotting proteins.
\synonyms
Kasabach-Merritt Phenomenon; Vascular Tumor Consumptive Coagulopathy

\medterm{Kasai Procedure} An operation for babies with biliary atresia, in which the blocked bile ducts outside the liver are removed and a loop of intestine is connected to the liver so bile can drain. Earlier surgery gives the best chance of restoring bile flow, but many children eventually need a liver transplant.

\synonyms
Hepatoportoenterostomy; Kasai Operation; Portoenterostomy

\medterm{Kashin-Beck Disease} A long-term disease of the bones and joints that can cause pain, stiffness, deformity, and poor growth. It is linked to low selenium and environmental factors in some regions.

\medterm{Katayama Fever} An acute systemic reaction to schistosome infection, usually occurring weeks after exposure and causing fever, rash, cough, abdominal symptoms, and eosinophilia.

\synonyms
Katayama Syndrome; Acute Bilharziasis

\medterm{Katz Index} A scale that assesses independence in six basic activities of daily living: bathing, dressing, toileting, transferring, continence, and feeding.

\medterm{Kawasaki Disease} An acute inflammatory disease of childhood that can affect the coronary arteries. It commonly causes persistent fever, rash, conjunctival redness, mouth or extremity changes, and swollen lymph nodes; diagnosis is clinical and incomplete forms occur.

\synonyms
Mucocutaneous Lymph Node Syndrome; KD

\medterm{Kawasaki Disease Diagnostic Criteria} The usual clinical pattern used to identify Kawasaki disease: fever for at least five days plus at least four of five findings—red eyes without pus, a varied skin rash, swollen neck lymph nodes, changes in the lips or mouth, and redness, swelling, or peeling of the hands or feet. Some children have an incomplete form.

\medterm{Kayser-Fleischer Rings} Golden-brown or greenish rings around the colored part of the eye caused by copper buildup. They strongly suggest Wilson disease when other findings fit and are seen with a special eye examination.

\medterm{K-Complex} A large brain-wave pattern seen during stage 2 sleep. It may occur on its own or after a sound or other sensory stimulus and is recorded during an electroencephalogram (EEG).

\medterm{Kearns-Sayre Syndrome} A rare inherited disorder affecting the mitochondria, the parts of cells that make energy. It usually begins in childhood or young adulthood and can cause drooping or limited eye movement, vision problems, heart-rhythm problems, and muscle weakness.

\medterm{Keel Breast} A chest-wall shape in which the breastbone pushes outward, giving the chest a keel-like appearance. It usually reflects pectus carinatum, which develops during growth and may affect appearance, breathing, or exercise tolerance.

\medterm{Kegel} A pelvic-floor muscle contraction exercise used to improve pelvic-floor strength and control of selected types of urinary or fecal incontinence.

\medterm{Kegel Exercise} Repeated tightening and relaxing of pelvic-floor muscles to improve strength and control, helping reduce urine leakage and support pelvic organs.

\medterm{Kehr Sign} Pain felt at the tip of the left shoulder when blood or irritation affects the underside of the diaphragm, the main breathing muscle. It can occur with irritation from a spleen injury but is not specific to it.

\medterm{Keloid} A raised, firm scar that grows beyond the edges of the original wound because the body makes too much scar tissue. It may itch or hurt and can develop after surgery, injury, burns, or piercings.

\synonyms
Keloid Scar; Cheloid

\medterm{Keloid and Hypertrophic Scar} Raised scars caused by excessive collagen deposition during wound healing. Hypertrophic scars remain within the original wound, whereas keloids extend beyond it; symptoms may include itch, pain, tightness, and cosmetic concern.

\medterm{Keloid Scar} A raised, firm scar that grows beyond the original skin injury because of excessive scar formation. It may itch, hurt, or recur after removal and is more common in some families and skin types.

\medterm{Kennedy Disease} An X-linked recessive, adult-onset lower motor neuron disorder caused by a CAG trinucleotide repeat expansion in the androgen receptor (AR) gene, presenting with slowly progressive proximal limb and bulbar muscle weakness, tongue and perioral fasciculations, muscle cramps, and signs of androgen insensitivity (gynecomastia and testicular atrophy).
\synonyms
Spinal and Bulbar Muscular Atrophy; SBMA

\medterm{Kennedy's Disease} A rare inherited motor-neuron disease, also called spinal and bulbar muscular atrophy, caused by an androgen-receptor gene change. It mainly affects men and causes slowly progressive muscle weakness, wasting, cramps, and sometimes swallowing or speech problems.

\medterm{Keratan Sulfate} A complex sugar found mainly in the clear front of the eye, cartilage, and brain. In some inherited disorders it cannot be broken down properly and builds up, damaging the eyes, bones, or other tissues.

\medterm{Keratan Sulphate} A sulfated glycosaminoglycan found especially in the cornea and cartilage; abnormal accumulation occurs in some mucopolysaccharidoses.

\medterm{Keratectomy} Removal or reshaping of part of the cornea, the clear front surface of the eye. It may be done to correct vision or remove abnormal tissue, and it can affect eyesight.

\medterm{Keratic Precipitates} Small groups of inflammatory cells that collect on the inner surface of the cornea, usually during inflammation inside the front of the eye. Their size and appearance can help describe the type of inflammation.

\medterm{Keratin} A strong structural protein forming much of the skin’s outer layer, hair, nails, and other protective tissues. It helps resist water loss, friction, chemicals, and physical injury.

\medterm{Keratinocyte} The main type of cell in the outer layer of skin. It produces keratin and helps form the skin barrier that limits water loss and protects against injury and infection.

\medterm{Keratinocytes} The predominant cells of the epidermis that produce keratin and form the skin’s protective barrier.

\medterm{Keratin Pearl} A round, layered collection of keratin seen under a microscope in some well-differentiated squamous cell skin cancers.

\synonyms
Squamous Pearl; Horn Pearl

\medterm{Keratitis} Inflammation or infection of the clear front surface of the eye. It can cause eye pain, redness, tearing, sensitivity to light, blurred vision, or a cloudy spot.

\synonyms
Corneal Ulceration; Corneal Inflammation

\medterm{Keratoacanthoma} A fast-growing, dome-shaped skin growth with a central plug, usually on sun-exposed skin. It may look like a type of skin cancer, so a clinician often removes or samples it. Treatment depends on its appearance, size, site, and the test results.

\synonyms
Well-Differentiated Squamous Cell Carcinoma, Keratoacanthoma Type; KA; Crateriform Epithelioma

\medterm{Keratocele} Severe thinning or ulceration of the cornea that allows its inner layers to bulge outward. It can lead to a hole in the cornea and permanent loss of vision.

\medterm{Keratoconjunctivitis} Inflammation of the cornea and the conjunctiva, the clear surface covering the white of the eye. Infection, allergy, dryness, immune disease, or chemical injury may cause redness, pain, discharge, light sensitivity, and blurred vision.

\medterm{Keratoconjunctivitis Sicca} Long-lasting dryness and inflammation of the eye’s front surface, commonly called dry-eye disease. It can cause burning, grittiness, redness, fluctuating vision, or light sensitivity because too few tears are made or tears evaporate too quickly.

\synonyms
Dry Eye Syndrome; Ocular Surface Disease

\medterm{Keratoconus} A condition in which the clear front surface of the eye becomes thin and bulges into a cone shape. It can cause blurred or distorted vision, glare, and sensitivity to light. Glasses, special contact lenses, or treatment to strengthen the cornea may help.

\medterm{Keratocyte Count} An estimate of the number of living support cells in the middle layer of the cornea. It can help assess corneal health, injury, disease, or changes after eye surgery, although results are interpreted with the complete eye examination.

\medterm{Keratoderma} Excessive thickening of the skin on the palms or soles. It may be inherited or develop because of skin disease, infection, medicines, or another illness and can cause painful cracks or difficulty walking.

\medterm{Keratoderma Blennorrhagica} Thickened, scaly, sometimes pustular skin lesions, usually on the soles or palms, associated with reactive arthritis. Treatment focuses on the underlying infection or inflammation and the skin symptoms.

\medterm{Keratoderma Blennorrhagicum} Thickened, scaly, or pus-filled skin lesions on the palms and soles associated with reactive arthritis. The lesions may spread to other areas and often appear after an infection affecting the urinary or digestive tract.

\medterm{Keratoderma Blenorrhagicum} Alternate spelling of keratoderma blennorrhagicum, a cutaneous manifestation of reactive
arthritis with waxy, crusty, vesiculopustular lesions on the palms and soles. It is a
reactive phenomenon typically occurring in the context of genitourinary or
gastrointestinal infection.

\medterm{Keratodermia Blennorrhagicum} Thickened, scaly, or wart-like skin lesions, often on the palms and soles, associated with reactive arthritis and sometimes infection-triggered inflammation.

\medterm{Keratolysis} Breakdown or loss of keratinized tissue, especially the superficial layer of the skin; it may occur in some infections and inflammatory skin disorders.

\medterm{Keratolytics} Substances, such as salicylic acid or urea, that soften and loosen excess hard skin. They may be used for thickened skin, corns, calluses, or some warts.

\medterm{Keratoma} A localized thickening or tumor-like growth of keratinized skin, such as a callus or corn; the exact type and need for evaluation depend on its appearance and behavior.

\medterm{Keratomalacia} Softening and ulceration of the cornea caused most often by severe vitamin A deficiency, potentially leading to blindness.

\medterm{Keratometer} An eye instrument that measures the curvature of the cornea, the clear front surface of the eye. The measurement helps assess astigmatism and fit rigid contact lenses.

\medterm{Keratometry} Measurement of the curve of the cornea, the clear front surface of the eye. It helps assess focusing problems and choose contact lenses or an artificial lens for cataract surgery.

\medterm{Keratoplasty} Surgery to replace or reshape part or all of the cornea, the clear front surface of the eye. It can restore vision after corneal damage, but healing problems, infection, and rejection of donor tissue are possible.

\medterm{Keratoprosthesis} An artificial cornea implanted to restore sight when a donor-corneal transplant is unlikely to work, such as after severe chemical injury or extensive surface damage. It carries risks of infection, glaucoma, and device problems.

\medterm{Keratosis} A thickened, rough, or otherwise abnormal area of skin caused by excess keratin. Some forms are harmless, while others result from infection, inflammation, sun damage, or a change that can precede skin cancer.

\medterm{Keratosis Follicularis} A condition in which extra keratin blocks hair follicles and causes many rough, small bumps. The name may describe keratosis pilaris or a rarer inherited skin disorder, so the appearance and history determine the diagnosis.

\medterm{Keratosis Follicularis (Darier Disease)} An inherited disorder of epidermal cell adhesion that causes persistent greasy, crusted, brown keratotic papules in seborrheic areas, nail changes, and sometimes mucosal lesions. Heat, sweating, friction, and ultraviolet exposure may worsen it.

\medterm{Keratosis Obturans} A buildup of shed skin cells that blocks the ear canal. It can cause ear pain, hearing loss, discharge, and widening or damage of the canal.

\medterm{Keratosis Pilaris} A common, harmless skin condition that causes small, rough bumps around hair follicles, usually on the upper arms, thighs, buttocks, or face.

\synonyms
KP

\medterm{Kerion} An inflamed, boggy, pus-forming mass caused by a fungal infection of hair-bearing skin, usually a severe form of tinea capitis.

\medterm{Kerley A Lines} Long, fine lines seen on a chest X-ray, usually spreading outward from the central lung areas. They can reflect fluid buildup in the lung tissue.
\synonyms
Kerley's A Lines; Upper Zone Interstitial Edema Lines

\medterm{Kerley B Lines} Short, thin lines seen near the outer edges and bases of the lungs on a chest X-ray. They commonly reflect fluid buildup in the lung tissue, often related to heart failure.
\synonyms
Kerley's B Lines; Septal Lines

\medterm{Kernicterus} Permanent brain damage caused by very high levels of bilirubin, a yellow substance that builds up during severe newborn jaundice. It can cause movement problems, hearing loss, problems looking upward, and damage to tooth enamel.

\synonyms
Bilirubin Encephalopathy; Nuclear Jaundice; Chronic Bilirubin Toxicity

\medterm{Kernig Sign} Pain or resistance when the knee is straightened while the hip is bent. It may suggest irritation of the coverings of the brain and spinal cord, such as with meningitis, but it is not diagnostic by itself.

\medterm{Kernig's Sign} Alternate spelling; see \textit{Kernig Sign}.

\medterm{Kernohan Notch Phenomenon} A paradoxical false-localizing neurological sign in which acute uncal transtentorial herniation displaces the midbrain against the contralateral rigid tentorial notch (Kernohan notch), compressing the contralateral cerebral peduncle (corticospinal tract) and producing an unexpected ipsilateral hemiparesis.
\synonyms
Kernohan-Woltman Notch Phenomenon; Kernohan Sign

\medterm{Kerosene Poisoning} Harm after swallowing or inhaling kerosene. The greatest danger is entry into the lungs, which can cause chemical pneumonia; coughing, choking, breathing difficulty, drowsiness, or vomiting may occur.

\medterm{Keshan Disease} A heart-muscle disease linked to severe selenium deficiency, occurring mainly in regions where soil and food contain little selenium. It can weaken the heart and cause heart failure.

\medterm{Ketamine} Ketamine is a dissociative anaesthetic and NMDA-receptor antagonist used for anaesthesia, procedural sedation, and selected treatment-resistant depression; emergence reactions, increased blood pressure, nausea, misuse, and respiratory complications are possible.

\medterm{Ketamine Hydrochloride} Ketamine hydrochloride is a dissociative anesthetic used for anesthesia, procedural sedation, and selected pain or depression treatments; it can cause increased blood pressure, hallucinations, and impaired coordination.

\medterm{Ketamine Toxicity} Acute adverse effects or overdose of the NMDA receptor antagonist dissociative anesthetic ketamine, characterized by perceptual distortions, hallucinations, severe emergence delirium, profound nystagmus, hypertension, tachycardia, and in chronic abusers, severe ulcerative cystitis ("ketamine bladder").

\medterm{Keto Acid} A chemical made when the body breaks down certain amino acids or other nutrients. Keto acids can be used to make energy or new substances, but abnormal amounts may occur in metabolic disease.

\medterm{Ketoacidosis} A dangerous buildup of acidic ketones in the blood, usually because there is too little insulin or prolonged starvation. It can cause vomiting, abdominal pain, deep breathing, confusion, coma, and death.

\medterm{Ketoconazole} Ketoconazole is an imidazole antifungal that inhibits fungal ergosterol synthesis; topical forms treat superficial mycoses, while oral use is restricted because of potentially severe liver injury, adrenal suppression, and drug interactions.

\medterm{Ketogenesis} The production of ketone bodies from fatty acids, mainly in the liver during fasting, prolonged exercise,
or inadequate insulin activity.

\medterm{Ketogenic Diet} An eating plan very low in carbohydrates and higher in fat that makes the body use ketones for energy. It can help some children with difficult-to-control epilepsy, but may cause nutrient deficiencies, constipation, kidney stones, or abnormal cholesterol without supervision.

\medterm{Ketone} A type of chemical compound. In medicine, the term often refers to ketone bodies made when the body uses fat for energy; small amounts can be normal, but high levels may cause ketoacidosis.

\medterm{Ketone Bodies} Water-soluble molecules produced mainly by the liver from fatty-acid-derived acetyl-CoA
during fasting, starvation, prolonged exercise, or insulin deficiency. The principal
ketone bodies are acetoacetate, beta-hydroxybutyrate, and acetone.

\medterm{Ketones} Water-soluble substances produced mainly by the liver when the body breaks down fat for energy, such as during fasting, a low-carbohydrate state, or insufficient insulin. The main ketone bodies are beta-hydroxybutyrate, acetoacetate, and acetone. Small amounts can be normal, but high levels may cause ketoacidosis, especially in diabetes; ketones can be measured in blood or urine.

\medterm{Ketones Blood Test} A blood test measuring ketone bodies, such as beta-hydroxybutyrate, used to assess ketosis and help diagnose or monitor diabetic ketoacidosis and other causes of ketone production.

\medterm{Ketones Urine Test} A urine test detecting ketone bodies produced when the body metabolizes fat; it can support evaluation of diabetic ketoacidosis, starvation, vomiting, or other metabolic states but may lag behind blood ketone levels.

\medterm{Ketonuria} Ketonuria is the presence of ketone bodies in urine, occurring with fasting, vomiting, low-carbohydrate intake, pregnancy, or uncontrolled diabetes; marked ketonuria with hyperglycaemia may indicate diabetic ketoacidosis.

\medterm{Ketoprofen} A non-steroidal anti-inflammatory drug that inhibits cyclo-oxygenase and reduces pain, inflammation, and fever; gastrointestinal bleeding, kidney injury, cardiovascular events, and hypersensitivity are important risks.

\medterm{Ketoprofen Overdose} Harm caused by taking too much ketoprofen, a nonsteroidal anti-inflammatory drug. It may cause nausea, vomiting, stomach bleeding, drowsiness, kidney injury, seizures, or a dangerous buildup of acid in the blood.

\medterm{Ketorolac} A potent non-steroidal anti-inflammatory analgesic used for short-term management of moderate-to-severe acute pain; treatment duration is restricted because gastrointestinal bleeding, kidney injury, and impaired healing can occur.

\medterm{Ketorolac Tromethamine} The tromethamine salt of ketorolac, a potent non-steroidal anti-inflammatory analgesic used short term for moderate-to-severe acute pain; it is contraindicated in several bleeding, renal, pregnancy, and perioperative situations.

\medterm{Ketosis} A state in which the body makes and uses increased amounts of ketones after breaking down fat for energy. It can occur during fasting or a low-carbohydrate diet, but high ketone levels with illness or high blood sugar may signal diabetic ketoacidosis.

\medterm{Ketosis-Prone Type 2 Diabetes} A form of diabetes that may present with diabetic ketoacidosis despite features of type 2 diabetes, such as obesity or preserved insulin production. Insulin dependence can change over time according to beta-cell reserve, autoimmunity, and recovery from the initial metabolic stress.
\medterm{Ketotic Hypoglycemia} Low blood sugar accompanied by ketones, usually after a young child has gone without food or has been ill and unable to eat. Drowsiness, weakness, vomiting, or seizures may occur. Repeated episodes can have other causes.

\medterm{Ketotifen} An antihistamine with mast-cell-stabilizing activity used in some countries for allergic disorders and, in ophthalmic form, allergic conjunctivitis; drowsiness, increased appetite, and irritation may occur.

\medterm{Kidney} One of two bean-shaped organs that filter waste from the blood into urine. The kidneys also help control body water, salts, blood pressure, and the balance of acids in the blood.

\medterm{Kidney Abscess} Alternate name for renal abscess, a localized collection of pus in or around the kidney.

\synonyms
Renal Abscess; Perinephric Abscess

\medterm{Kidney Biopsy} A procedure in which a small sample of kidney tissue is removed, usually with a needle guided by ultrasound, to identify kidney disease. Bleeding is the main risk.

\medterm{Kidney Calculi} Alternate term; see \textit{Kidney stones}.

\medterm{Kidney Cancer} A malignant tumor arising in the kidney. Renal cell carcinoma is the most common type in adults, but other types occur, especially in children.

\medterm{Kidney Cancer In Children} Cancer arising in a child’s kidney. Wilms tumor is the most common type, but other kidney tumors can occur; symptoms may include an abdominal mass, pain, blood in the urine, fever, or high blood pressure.

\medterm{Kidney Cysts} Fluid-filled sacs in or on the kidneys. Simple cysts are common and usually harmless, while complex or multiple cysts may be associated with infection, bleeding, inherited disease, or cancer risk and may impair kidney function.

\synonyms
Renal Cysts; Renal Cystic Disease

\medterm{Kidney Disease} Any condition that damages the kidneys or reduces their ability to filter blood, balance fluids and minerals, control blood pressure, or make hormones. Causes include diabetes, high blood pressure, infection, inflammation, inherited disorders, and medicines.

\medterm{Kidney Disease and Pregnancy} Kidney disease can affect fertility, blood pressure, fluid balance, medication choices, and pregnancy outcomes, while pregnancy can worsen some kidney disorders. Preconception counseling and coordinated obstetric-nephrology care help monitor kidney function, proteinuria, blood pressure, fetal growth, and medicine safety.

\medterm{Kidney Disease-Related Anemia} Anemia caused or worsened by kidney disease, mainly because damaged kidneys produce less erythropoietin, the hormone that stimulates red-cell production. Iron deficiency, inflammation, and blood loss may also contribute.

\medterm{Kidney Diseases} Disorders affecting the kidneys or urinary collecting system that can impair filtration, urine formation, fluid balance, blood pressure, or electrolyte regulation.

\medterm{Kidney Failure} A condition in which the kidneys can no longer remove enough waste or control fluid and body salts. It may develop suddenly or gradually and can require medicines, dialysis, or a transplant.

\synonyms
End-Stage Renal Disease; ESRD; Renal Failure

\medterm{Kidney Function Tests} Blood and urine tests used to assess how well the kidneys remove waste and control fluid and salts. Common measures include creatinine, estimated filtration rate, urine protein, and urinalysis.

\medterm{Kidney Infection} An infection of one or both kidneys, usually caused by bacteria that travel upward from the bladder. It can cause fever, chills, side or back pain, nausea, vomiting, and painful urination and may lead to sepsis.

\medterm{Kidney Neoplasm} An abnormal growth in the kidney that may be noncancerous or cancerous. It may cause blood in the urine, pain, a lump, weight loss, or no symptoms; imaging and sometimes biopsy identify its type and extent.

\medterm{Kidney Pain} Pain felt in the side or back below the ribs that may be caused by kidney stones, infection, obstruction, inflammation, or referred pain.

\synonyms
Flank Pain; Renal Colic; Nephralgia

\medterm{Kidney Position} A side-lying surgical position in which the waist is gently flexed and the side of the body is raised. It gives surgeons better access to the kidney and nearby tissues while staff protect pressure points and breathing.

\medterm{Kidney Removal} Surgery to remove part or all of a kidney, usually for cancer, severe damage, infection, or a nonfunctioning kidney. The remaining kidney may provide adequate function, but bleeding, infection, injury to nearby organs, or reduced kidney function can occur.

\medterm{Kidney Replacement Therapy} Treatment that replaces some functions of severely damaged kidneys. It includes hemodialysis, peritoneal dialysis, continuous kidney replacement therapy, and kidney transplantation.

\medterm{Kidney Stone} Singular form; see \textit{Kidney stones}.

\medterm{Kidney Stones} Hard mineral deposits formed from substances in urine. They may obstruct the urinary tract and cause colicky flank or abdominal pain, blood in the urine, vomiting, or infection.

\medterm{Kidney Stones In Children} Mineral deposits that form in a child’s kidney or urinary tract. They may cause side or abdominal pain, blood in the urine, nausea, vomiting, urinary symptoms, or infection, but some cause no symptoms.

\synonyms
Pediatric Nephrolithiasis; Pediatric Urolithiasis; Pediatric Renal Calculi

\medterm{Kidney Transplant} Surgical placement of a healthy donor kidney into a person with kidney failure to restore renal function. Lifelong immunosuppressive treatment is generally required to reduce rejection risk.

\medterm{Kidney Transplantation} Placement of a donor kidney to replace the filtering work of kidneys that have permanently failed. A living-donor kidney may work sooner, but all recipients need lifelong medicines and monitoring for rejection, infection, blood-vessel problems, and return of kidney disease.

\medterm{Kidney Transplant Rejection} An immune attack that damages a transplanted kidney. It may occur suddenly or gradually and can cause reduced urine, swelling, high blood pressure, or worsening blood tests; treatment depends on the type and severity.

\medterm{Kidney Trauma} Alternate name for renal trauma, an injury to the kidney caused by blunt or penetrating force.

\medterm{Kienbock Disease} Damage and gradual collapse of the lunate, a small bone in the wrist, because its blood supply is reduced. It may cause wrist pain, swelling, stiffness, and weaker grip. Treatment depends on bone damage, wrist alignment, symptoms, and whether arthritis has developed.

\medterm{Kieselguhr} Diatomaceous earth containing substantial silica; inhalation of its dust can produce occupational pneumoconiosis and other lung injury.

\medterm{Kiesselbach Plexus} A network of small arteries on the front part of the nasal septum and a common source of nosebleeds.

\medterm{Kiesselbach Triangle} An anterior vascular network on the anteroinferior nasal septum (Little's area) formed by anastomoses of the anterior ethmoidal, sphenopalatine, greater palatine, and superior labial arteries, serving as the anatomical site for over 90% of anterior epistaxis episodes.
\synonyms
Little Area; Kiesselbach's Plexus

\medterm{Kikuchi Disease} A usually short-term inflammation of the lymph nodes, most often in the neck, with no known cause. It mainly affects young adults and may cause tender swollen nodes, fever, night sweats, and a low white-blood-cell count.

\synonyms
Kikuchi-Fujimoto Disease; Histiocytic Necrotizing Lymphadenitis

\medterm{Kikuchi-Fujimoto Disease} Another name for Kikuchi disease, a usually short-term inflammation of the lymph nodes, most often in the neck. It may cause painful swollen nodes, fever, night sweats, and a low white-blood-cell count.
\synonyms
Kikuchi Disease; Histiocytic Necrotizing Lymphadenitis

\medterm{Killer Cold Virus} Alternate name; see \textit{Adenovirus Infection}.

\medterm{Killian's Triangle} An area of potential weakness in the pharyngeal wall between the thyropharyngeus and
cricopharyngeus muscles, predisposing to Zenker diverticulum formation.

\medterm{Killian Triangle} A triangular anatomical area of muscular weakness in the posterior hypopharyngeal wall situated between the oblique fibers of the thyropharyngeus muscle and the transverse fibers of the cricopharyngeus muscle; the site of mucosal herniation forming a Zenker diverticulum.

\synonyms
Killian Dehiscence

\medterm{Killip Classification} A four-level system for describing heart failure during a heart attack. The levels range from no signs of heart failure to lung fluid buildup and, at the most severe level, shock caused by the heart’s inability to pump enough blood.

\medterm{Kilocalorie} A unit of energy equal to 1,000 small calories; the dietary Calorie, written with a capital C, is one kilocalorie.

\medterm{Kilojoule} The SI unit of energy, equal to 1,000 joules; approximately 4.184 kilojoules equal one kilocalorie.

\medterm{Kimmelstiel-Wilson Lesion} A round area of thickened tissue in a kidney filter that can occur with long-standing diabetes. It is a sign of diabetic kidney damage and may occur with protein in the urine or reduced kidney function.
\synonyms
Nodular Diabetic Glomerulosclerosis; Kimmelstiel-Wilson Nodules

\medterm{Kimmelstiel-Wilson Nodules} Round areas of thickened tissue seen in the kidney filters of some people with long-standing diabetes. They reflect diabetic kidney damage and may occur with protein in the urine and declining kidney function.

\medterm{Kimmelstiel-Wilson Syndrome} A pattern of kidney scarring caused by long-standing diabetes. It can make protein leak into the urine and gradually reduce kidney function, sometimes leading to kidney failure.

\medterm{Kimura Disease} A rare long-term immune disorder causing painless lumps under the skin, usually in the head or neck, enlarged nearby lymph nodes, and high levels of eosinophils or IgE. It can sometimes affect the kidneys.

\medterm{Kinase} An enzyme that attaches a small chemical group called a phosphate to another molecule. This acts like a switch controlling cell growth, energy use, movement, communication, and many other processes.

\medterm{Kinase Inhibitor} A targeted therapeutic agent (typically a small molecule or monoclonal antibody) that selectively inhibits the enzymatic phosphorylation activity of specific protein kinases (such as tyrosine kinases, serine/threonine kinases, or cyclin-dependent kinases) involved in oncogenic cellular signaling cascades.

\medterm{Kinesia Paradoxa} A brief, unexpected ability to move normally or quickly in a person who usually has severe slowness or difficulty moving, as can occur in advanced Parkinson disease. It may happen during a strong emotional experience or in response to an urgent situation.
\synonyms
Paradoxical Kinesia; Paradoxical Movement

\medterm{Kinesthesia} The sense of how the body's parts are moving and where they are positioned, even when the eyes are closed. Information from muscles, joints, and skin helps the brain control posture, balance, and coordinated movement.

\medterm{Kinetics} The study of how a process changes over time. In medicine, pharmacokinetics describes how the body absorbs, distributes, changes, and removes a medicine, helping determine its dose and timing.

\medterm{Kinetochore} A specialized, disc-shaped multiprotein complex assembled onto the centromeric DNA region of a eukaryotic chromosome that captures spindle microtubules and coordinates chromosome segregation during mitosis and meiosis.

\medterm{Kingella} Kingella is a genus of bacteria that can live in the mouth and can occasionally cause bloodstream, bone, joint, or heart-valve infections.

\medterm{Kingella Denitrificans} Kingella denitrificans is a rare bacterium found mainly in the upper airway; it can occasionally cause invasive infection.

\medterm{Kingella Kingae} Kingella kingae is a bacterium that can cause bone or joint infection, especially in young children, and can rarely cause endocarditis.

\medterm{Kingella Oralis} Kingella oralis is a mouth-associated bacterium that rarely causes invasive infection, including bloodstream infection or endocarditis.

\medterm{Kingella Potus} Kingella potus is a rare bacterium that has been associated with occasional human infections; clinical importance depends on the specimen and evidence of true infection.

\medterm{King's College Criteria} A set of findings used to estimate the risk of death in acute liver failure and the possible need for a liver transplant. The criteria differ for acetaminophen poisoning and other causes.

\medterm{Kinin} A short peptide, such as bradykinin, that widens blood vessels, increases vascular permeability, and contributes to pain and inflammation.

\medterm{Kinin-Kallikrein System} An endogenous enzymatic cascade of blood proteins that produces bioactive peptide kinins (primarily bradykinin and kallidin) from high- and low-molecular-weight kininogens, mediating vasodilation, increased microvascular permeability, tissue inflammation, and pain sensation.

\synonyms
Kallikrein-Kinin Cascade

\medterm{Kinins} A group of vasoactive peptides, chiefly bradykinin and kallidin, generated from kininogens during inflammation; they promote vasodilation, increased vascular permeability, smooth-muscle effects, and pain.

\medterm{Kinocilium} A specialized hair-like structure on certain sensory cells, especially in the inner ear. It helps the cell detect the direction of movement and contributes to balance and hearing signals.

\medterm{Kissing Disease} Informal term; see \textit{Infectious mononucleosis}.

\medterm{Kite Angle} An anatomical radiographic angle measured on an anteroposterior radiograph of the foot between the longitudinal axis of the talus and the longitudinal axis of the calcaneus. A normal angle ranges from 20 to 40 degrees; convergence or reduction of this angle below 15 degrees signifies hindfoot parallelism, a diagnostic hallmark of congenital talipes equinovarus (clubfoot).

\synonyms
Talocalcaneal Angle; AP Kite Angle

\medterm{Klatskin Tumor} A cancer arising where the right and left bile ducts join near the liver. It can block bile flow and cause jaundice, itching, pale stools, or weight loss.

\medterm{Klebsiella} A genus of Gram-negative bacteria that can cause pneumonia, urinary-tract infection, bloodstream infection, and other healthcare-associated infections.

\medterm{Klebsiella Oxytoca} Klebsiella oxytoca is a bacterium that can cause urinary, bloodstream, wound, and hospital-acquired infections and can sometimes cause antibiotic-associated hemorrhagic colitis.

\medterm{Klebsiella Pneumonia} Pneumonia caused by Klebsiella bacteria. It can cause fever, cough, chest pain, breathlessness, and sometimes thick or blood-stained sputum; severe infection may destroy lung tissue and is more likely in people with serious illness or weakened defenses.

\medterm{Klebsiella Pneumoniae} Klebsiella pneumoniae is a bacterium that can cause pneumonia, urinary infection, bloodstream infection, and liver abscess; some strains are highly antibiotic-resistant.

\medterm{Klebsiella Rhinoscleromatis} A gram-negative bacillus causing chronic granulomatous infection of the nasal passages
and upper airway, leading to rhinoscleroma with progressive nasal obstruction and tissue
destruction.

\medterm{Klebsiella Singaporensis} Klebsiella singaporensis is a rare environmental bacterium that has occasionally been recovered from human clinical specimens; its significance depends on the specimen and symptoms.

\medterm{Klebsiella Variicola} Klebsiella variicola is a bacterium closely related to K. pneumoniae that can cause human infections, including bloodstream and urinary infections.

\medterm{Kleihauer-Betke Test} A quantitative acid-elution blood test used to detect and calculate the volume of fetomaternal hemorrhage; exposure to an acid buffer elutes adult hemoglobin (HbA) from maternal erythrocytes while acid-resistant fetal hemoglobin (HbF) inside fetal red blood cells remains intact and stains dark pink, allowing calculation of necessary Rh immune globulin (RhoGAM) dosing.
\synonyms
Acid Elution Test; Fetomaternal Hemorrhage Assay

\medterm{Kleine-Levin Syndrome} A rare disorder causing repeated episodes of excessive sleep, often with confusion, reduced activity, and changes in behavior or appetite. Episodes last days to weeks and are separated by periods of more typical function.

\medterm{Kleptomania} A mental-health disorder involving recurrent inability to resist impulses to steal items that are not needed for personal use or monetary value. The acts are typically preceded by rising tension and followed by relief or gratification, and cause distress or impairment.

\synonyms
Compulsive Stealing
Pathological Stealing

\medterm{Klinefelter's Syndrome} Alternate name; see \textit{Klinefelter Syndrome}.

\medterm{Klinefelter Syndrome} A genetic condition in males caused by one or more extra X chromosomes, most often the 47,XXY pattern. It can cause small testes, low testosterone, infertility, tall stature, less facial and body hair, breast enlargement, weak bones, and learning or language difficulties. A chromosome test identifies the condition.
\synonyms{47,XXY Syndrome; XXY Syndrome; Klinefelter-Reifenstein-Albright Syndrome}

\medterm{Klippel-Feil Syndrome} A condition present from birth in which two or more neck bones are joined together. It may cause a short neck, low hairline, limited neck movement, hearing loss, or spinal and kidney problems.

\medterm{Klippel-Trenaunay Syndrome} A rare congenital vascular disorder characterized by a classic triad of features: capillary malformations (port-wine stains), soft tissue and bone enlargement (hypertrophy, usually affecting one leg or arm), and slow-flow venous and lymphatic malformations (such as persistent embryonic veins, prominent varicose veins, and lymphedema). It is typically caused by non-inherited (somatic) mutations in the \textit{PIK3CA} gene. Complications include limb length discrepancy, chronic pain, recurrent skin infections (cellulitis), and blood clots (deep vein thrombosis and pulmonary embolism). Management is multidisciplinary and includes compression garments, physiotherapy, vascular procedures, and orthopedic care.

\synonyms
KTS; Klippel-Trénaunay Syndrome; Angio-Osteohypertrophy Syndrome; KTW Syndrome; KTWS

\medterm{Klippel-Trenaunay-Weber Syndrome} A congenital vascular and limb-overgrowth disorder with capillary malformation, venous or lymphatic abnormalities, and enlargement of a limb or its digits. It can cause pain, bleeding, ulcers, clotting, or pulmonary embolism and requires vascular assessment.

\medterm{Kl\"{u}ver-Bucy Syndrome} A profound neurobehavioral syndrome resulting from bilateral damage or surgical ablation of the anterior temporal lobes, specifically including the amygdaloid nuclei and hippocampus. Clinical hallmarks include visual agnosia (psychic blindness), compulsive oral exploration of objects (hyperorality), hypermetamorphosis, marked passivity or docility, altered dietary preferences, and inappropriate hypersexuality.

\synonyms
Kluver-Bucy Syndrome; Bilateral Amygdala Ablation Syndrome

\medterm{Klumpke Palsy} An injury to the lower trunk of the brachial plexus involving spinal nerve roots C8 and T1, typically resulting from upward traction on an abducted arm during birth or trauma, causing paralysis of intrinsic hand muscles resulting in a 'claw hand' deformity and ipsilateral Horner syndrome (if sympathetic T1 root fibers are disrupted).
\synonyms
Klumpke-Dejerine Palsy; Lower Trunk Brachial Plexus Injury

\medterm{Klumpke Paralysis} A nerve injury affecting the lower part of the network that supplies the arm and hand. It can cause severe hand weakness, numbness along the inner forearm and hand, and occasionally drooping of the eyelid and a smaller pupil on one side.

\medterm{Klumpke's Paralysis} Lower brachial-plexus palsy involving mainly the C8 and T1 nerve roots, causing hand weakness and sensory loss; Horner syndrome may accompany severe cases.

\medterm{Kluyvera} A genus of gram-negative bacteria in the Enterobacterales group; species can rarely cause opportunistic human infections.

\medterm{Kluyvera Ascorbata} A gram-negative Kluyvera species found in the environment and occasionally associated with urinary, bloodstream, or other opportunistic infections.

\medterm{Kluyvera Cryocrescens} A gram-negative environmental bacterium that rarely causes opportunistic human infection.

\medterm{Kluyvera Georgiana} A rarely encountered Gram-negative Kluyvera species that may appear in clinical samples. It can occasionally cause opportunistic infection, especially in people with serious illness, weakened defenses, or invasive medical devices.

\medterm{Kluyvera Intermedia} A Gram-negative bacterium occasionally found in human samples and rarely causing opportunistic urinary, wound, bloodstream, or other infections.

\medterm{Knee} The synovial joint between the femur, tibia, and patella, stabilized by the cruciate and collateral ligaments and menisci; it permits flexion and extension and bears body weight.

\medterm{Knee Arthroplasty} Surgery that replaces damaged parts of the knee joint with artificial components. It is usually performed for severe pain and loss of function from arthritis or major joint damage.

\medterm{Knee Arthroscopy} Knee arthroscopy is a minimally invasive procedure in which a camera and small instruments are inserted into the knee to diagnose or treat selected joint problems.

\medterm{Knee Bursitis} Inflammation of a fluid-filled sac around the knee, causing local swelling, tenderness, and pain with kneeling or movement. Repeated pressure, injury, arthritis, gout, or infection can contribute; fever or redness may occur with infection.

\synonyms
Prepatellar Bursitis; Housemaid's Knee; Pes Anserine Bursitis

\medterm{Kneecap} A flat, rounded bone at the front of the knee, embedded in the tendon of the thigh muscles. It protects the joint and improves the leverage used to straighten the knee; injury or arthritis can cause pain and restricted movement.

\medterm{Kneecap Dislocation} An injury in which the patella slips completely out of the groove at the end of the thighbone, usually toward the outside. It causes sudden pain, deformity, swelling, and difficulty walking and needs prompt assessment.

\medterm{Knee-Chest Position} A position in which the patient rests on the knees and chest, with the hips flexed; it may be used for spinal or anorectal procedures.

\medterm{Knee CT Scan} A detailed X-ray scan of the knee made by a computer. It is useful for complex fractures, bone alignment, arthritis, infection, masses, and selected soft-tissue problems, but it exposes the person to more radiation than a plain X-ray.

\medterm{Knee Dislocation} Complete anatomical disruption and loss of articulation between the distal femur and proximal tibia, representing a severe orthopaedic limb-threatening emergency frequently associated with multi-ligament rupture and high incidence of popliteal artery injury and peroneal nerve palsy.

\medterm{Knee Disorders} Problems affecting the knee joint, including injuries to ligaments or menisci, arthritis, bursitis, tendon disorders, infection, and fractures. They may cause pain, swelling, stiffness, instability, locking, or difficulty bearing weight.

\medterm{Knee Effusion} The abnormal accumulation of excess synovial fluid, pus, or blood within the joint capsule of the knee, clinically detected by the ballottement of the patella or bulge sign, warranting diagnostic arthrocentesis to differentiate inflammatory, septic, or hemarthrotic causes.

\medterm{Knee Hyperextension} Bending the knee backward beyond its normal range. It can strain or tear ligaments and may cause pain, swelling, instability, or difficulty bearing weight.

\medterm{Knee Injury And Meniscus Tears} A meniscus tear is damage to the cartilage cushion between the thighbone and shinbone, often after twisting or with degeneration. Pain, swelling, catching, or locking may occur; treatment ranges from rehabilitation to surgery according to tear, symptoms, and stability.

\medterm{Knee-Jerk} The automatic kick of the lower leg when the tendon below the kneecap is tapped. It tests part of the nerve pathway from the thigh and lower spinal cord and can help identify nerve or spinal problems.

\medterm{Knee Jerk Reflex} A monosynaptic deep tendon stretch reflex elicited by percussing the patellar ligament below the patella with a reflex hammer, stretching muscle spindles within the quadriceps femoris and provoking reflex knee extension mediated through the femoral nerve and spinal segments L2, L3, and L4.
\synonyms
Patellar Reflex; Quadriceps Reflex

\medterm{Knee Joint} A synovial joint formed mainly by the femur, tibia, and patella. It permits flexion and extension with limited rotation and is
stabilized by ligaments, menisci, muscles, and the joint capsule.

\medterm{Knee Joint Replacement} Surgery replacing damaged surfaces of the knee with artificial parts to relieve severe pain and improve movement. It is commonly used for advanced arthritis; infection, blood clots, stiffness, loosening, persistent pain, or repeat surgery are possible.

\medterm{Knee Meniscus} A pair of crescent-shaped fibrocartilaginous pads (medial and lateral menisci) interposed between the femoral condyles and the tibial plateau, functioning to distribute compressive axial loads, absorb shock, lubricate the knee joint, and deepen the tibial articular surfaces.

\medterm{Knee Microfracture Surgery} An operation that makes small holes in the bone beneath a damaged area of knee cartilage so marrow cells can form repair tissue. It may improve symptoms in selected small defects, but the new tissue is not identical to normal cartilage and results vary.

\medterm{Knee MRI Scan} Magnetic resonance imaging of the knee using a magnetic field and radio waves to evaluate cartilage, menisci, ligaments, tendons, bone marrow, muscles, and joint fluid.

\medterm{Knee Pain} Pain felt in or around the knee, commonly caused by injury, osteoarthritis, bursitis, tendinitis, or inflammatory disease. Swelling, instability, inability to bear weight, or fever may also occur.

\synonyms
Gonalgia; Knee Arthralgia

\medterm{Knee Prosthesis} An artificial joint component used to replace damaged surfaces of the knee, usually including femoral and tibial components and sometimes a patellar component.

\medterm{Knee Replacement} An operation that replaces part or all of a damaged knee joint with an artificial joint to reduce pain and improve movement. Recovery includes wound care, exercises, and rehabilitation.

\medterm{Kniest Dysplasia} A rare inherited disorder of bone and cartilage growth that causes short stature, enlarged joints, unusual facial features, and progressive joint problems. Some people also develop serious eye, hearing, or breathing complications.

\medterm{Knife-Cut Ulcers From Crohn’S Disease} Deep, narrow ulcers in the intestinal lining that look like cuts and can occur in Crohn disease. They may cause bleeding, narrowing of the bowel, abnormal connections, pain, or blockage and are identified during intestinal examination.

\medterm{Knock-Knee} Genu valgum, a leg alignment in which the knees angle inward and may touch while the ankles remain apart.

\medterm{Knock Knees} Medial angulation of the knees causing the knees to touch while the ankles remain apart,
also called genu valgum. It is common during childhood development but may be pathologic
when severe, asymmetric, persistent, or associated with pain or growth disturbance.

\synonyms
Genu Valgum

\medterm{Knotless Anchor} A small surgical anchor securing sutures in bone without requiring tied knots. Orthopedic surgeons use it to repair tendons, ligaments, or other soft tissues, with risks including loosening, infection, or tissue failure.

\medterm{Knuckle} A projecting finger joint, usually a metacarpophalangeal joint; it may be affected by trauma, inflammation, or arthritis.

\medterm{Knuckle Pad} A firm, usually painless thickening of skin and tissue over a finger joint. It is commonly related to repeated pressure or a fibrous tissue change and is generally harmless.

\medterm{Kocher Clamp} A surgical instrument with teeth used to hold firm tissue or clamp blood vessels. Its strong grip is useful during operations, but careless use can crush or injure delicate tissue.

\medterm{Kocher-Debre-Semelaigne Syndrome} A rare pediatric clinical syndrome associated with chronic, untreated congenital or juvenile hypothyroidism, characterized by generalized muscular pseudohypertrophy ('infant Hercules' phenotype), muscle stiffness, sluggish deep tendon reflexes, and myxedema.
\synonyms
Hypothyroid Muscular Pseudohypertrophy; Debré-Sémélaigne Syndrome

\medterm{Kocher-Debré-SéMéLaigne Syndrome} A rare condition in infants or young children with severe, untreated underactive thyroid disease. Muscles may look unusually large but are weak and slow, with delayed growth and development, puffy facial features, and other signs of hypothyroidism. Thyroid-hormone treatment usually improves the muscle changes and weakness.

\synonyms
Cretinism with Muscular Hypertrophy; Infantile Hypothyroid Myopathy

\medterm{Kocher Incision} An open surgical subcostal oblique abdominal incision made parallel to and approximately 2 to 3 cm below the right costal margin, traditionally utilized for open cholecystectomy and biliary tract surgical access.

\medterm{Kocher Maneuver} A surgical technique in which the peritoneum along the lateral border of the duodenum is sharply incised, mobilizing the duodenum and head of the pancreas medially and anteriorly to gain access to retroperitoneal structures including the inferior vena cava, aorta, and posterior pancreas.
\synonyms
Duodenal Mobilization; Kocherization

\medterm{Koch's Bacillus} The bacterium that causes tuberculosis, an infection most often affecting the lungs but capable of involving many organs.

\synonyms
Tubercle Bacillus; M. tuberculosis

\medterm{Koebnerisation} Development of lesions of an existing skin disease at sites of trauma or injury, classically seen in psoriasis, vitiligo, and lichen planus.

\medterm{Koebner Phenomenon} Development of lesions of a skin disease, classically psoriasis or vitiligo, at sites of injury or other trauma.

\medterm{Kohler Disease} A rare, self-limiting childhood osteochondrosis characterized by avascular necrosis of the tarsal navicular bone, occurring predominantly in boys aged 3 to 7 years, presenting with antalgic limp, medial midfoot pain, and localized dorsal swelling.

\synonyms
Kohler Bone Disease

\medterm{KOH Preparation} A test in which a skin, hair, or nail sample is treated with potassium hydroxide and viewed under a microscope to look for fungi. A negative result does not always rule out infection if the sample is small or treatment has already begun.

\medterm{Koilocyte} A cell with changes often seen after infection with human papillomavirus (HPV). Under a microscope, it has an enlarged or irregular nucleus and a clear space around the nucleus. The finding supports HPV-related change but is not enough by itself to determine the diagnosis.

\medterm{Koilonychia} Thin, brittle nails that become concave and curve upward at the edges, giving them a spoon-like shape. It is often linked to iron deficiency but can also result from poor nutrition, repeated injury, inherited conditions, or other illness; new or widespread changes need evaluation.

\medterm{Kombucha} A fermented tea made with a culture of bacteria and yeast. It may contain sugar, acids, caffeine, and a small amount of alcohol; contaminated or excessive products can cause illness, especially in vulnerable people.

\medterm{Konno Procedure} Open-heart surgery to widen a severely narrowed path through which blood leaves the left ventricle. The surgeon enlarges the area below and around the aortic valve, often replacing or repairing the valve at the same time. It is used for complex or recurrent obstruction and carries the usual risks of major heart surgery.

\synonyms
Konno-Rastan Procedure; Aortoventriculoplasty

\medterm{Koplik Spots} Small white or bluish spots on a red background inside the cheek, usually opposite the back teeth. They are characteristic of measles and often appear shortly before the skin rash.

\medterm{Koplik's Spots} Small white or bluish-gray spots on the inner cheek, often with a red base, that are characteristic of measles and may appear shortly before the skin rash.

\medterm{Korotkoff Sounds} Sounds heard through a stethoscope over an artery as a blood-pressure cuff is deflated; their appearance and disappearance are used to estimate systolic and diastolic pressure.

\medterm{Korsakoff's Syndrome} A chronic amnestic disorder, commonly associated with severe thiamine deficiency and alcohol-use disorder, causing impaired new learning, memory loss, and sometimes confabulation.

\medterm{Korsakoff Syndrome} A chronic amnestic disorder caused most often by severe thiamine deficiency associated with alcohol-use disorder. It causes impaired formation of new memories, loss of past memories, and sometimes confabulation.

\medterm{Kostmann Syndrome} An autosomal recessive form of severe congenital neutropenia caused by mutations in the HAX1 or ELANE genes, characterized by early arrest of granulopoiesis at the promyelocyte stage in the bone marrow, profound absolute neutropenia (less than 200/uL), and life-threatening bacterial infections starting in early infancy.
\synonyms
Severe Congenital Neutropenia Type 1; Infantile Genetic Agranulocytosis

\medterm{Koumpounophobia} An intense fear or strong distress related to buttons. It is a specific phobia when it causes significant anxiety or interferes with daily life, and it can be treated with psychological therapies such as exposure-based therapy.

\medterm{Kounis Syndrome} A heart attack or temporary narrowing of a heart artery triggered by a severe allergic reaction. It may cause chest pain together with rash, swelling, wheezing, or low blood pressure.

\medterm{KP} Abbreviation; see \textit{Keratosis Pilaris}.

\medterm{Krabbe Disease} A rare inherited disorder in which missing enzymes damage myelin, the protective covering of nerves. Early disease can cause irritability, stiffness, seizures, loss of skills, and progressive problems with movement, vision, or development.

\medterm{Krabbe's Disease} A rare inherited disorder in which missing enzymes cause myelin loss, progressively damaging nerves and producing weakness, stiffness, vision problems, or developmental decline.

\medterm{KRAS Gene} A gene that helps control cell growth and division. Certain changes can keep its growth signal switched on, contributing to cancers and affecting which targeted treatments are likely to work.

\medterm{Kraurosis} An older term for atrophic and sclerotic vulvar disease, often referring to lichen sclerosus and associated with itching, discomfort, and tissue fragility.

\medterm{Kraurosis Vulvae} Advanced progressive atrophy, shrinkage, and sclerosis of the female external genitalia occurring in chronic untreated vulvar lichen sclerosus, characterized by resorption of the labia minora, narrowing of the vaginal introitus, and thin parchment-like epidermal fragility.
\synonyms
Vulvar Lichen Sclerosus; Atrophic Vulvitis

\medterm{Krebs Cycle} A series of chemical reactions in mitochondria that oxidizes acetyl-CoA and generates energy-rich molecules used for cellular energy production.

\synonyms
Tricarboxylic Acid Cycle; TCA Cycle

\medterm{Kringles} Small folded parts of certain proteins that help them attach to other molecules. They are found in proteins involved in blood clot breakdown and in some growth signals.

\medterm{Krukenberg Spindle} A vertical, brownish line of pigment on the inner surface of the cornea. It may occur with pigment dispersion in the eye and can be associated with changes in eye pressure.
\synonyms
Endothelial Pigment Spindle; Pigment Dispersion Corneal Spindle

\medterm{Krukenberg's Tumor} Alternate spelling; see \textit{Krukenberg Tumor}.

\medterm{Krukenberg Tumor} An ovarian tumor caused by cancer that has spread from another organ, most often the stomach or digestive tract. It may cause abdominal swelling, pain, menstrual changes, or digestive symptoms.

\medterm{Kt/V} A number used to estimate how much urea a dialysis treatment removes. It combines the dialyzer's cleaning ability and treatment time with the person's body-water volume; clinicians use it with symptoms, fluid balance, laboratory results, and the dialysis schedule.

\medterm{Kugelberg-Welander Disease} A milder form of inherited spinal muscular atrophy, also called type 3 spinal muscular atrophy. It usually begins after 18 months of age and causes slowly worsening weakness in muscles close to the body, especially the hips and thighs. Children may have a waddling walk or trouble climbing stairs. Many walk independently, and life expectancy is often near normal with current care.

\synonyms
Spinal Muscular Atrophy Type 3 (SMA 3); Juvenile Spinal Muscular Atrophy

\medterm{Kulchitsky Cells} Specialized hormone-producing cells found mainly in the digestive and respiratory systems. They release chemical signals, including serotonin, that help regulate movement, secretion, and other organ functions.

\synonyms
EC Cells; Argentaffin Cells

\medterm{Kupffer Cells} Immune cells lining the small blood channels in the liver. They remove old blood cells, particles, and some germs from blood coming from the digestive organs.

\synonyms
Stellate Macrophages of the Liver

\medterm{Kupffer's Cell} A Kupffer cell is a resident macrophage lining the liver sinusoids; it removes microbes, damaged blood cells, and particles from portal blood and contributes to hepatic inflammation and iron handling.

\medterm{Kuru} Kuru is a rare fatal prion disease historically associated with ritual cannibalism, causing progressive ataxia, tremor, and neurologic decline.

\medterm{Kussmaul Breathing} Deep, rapid, labored breathing that helps remove extra carbon dioxide when the blood is dangerously acidic. It can occur in diabetic ketoacidosis, severe kidney failure, or other illnesses.

\medterm{Kussmaul Respiration} Alternate name; see \textit{Kussmaul Breathing}.

\medterm{Kussmaul Sign} A rise or failure to fall in the pressure of the large neck veins during breathing in. It suggests that the right side of the heart cannot fill normally and may occur with severe heart failure, constrictive pericarditis, or other conditions affecting right-heart filling.

\medterm{Kveim-Siltzbach Test} An old skin test once used to support the diagnosis of sarcoidosis. Material from sarcoid tissue was injected under the skin and the area was later biopsied. The test is rarely used now because of safety, availability, and accuracy concerns.
\synonyms
Kveim Test; Sarcoidosis Skin Test

\medterm{Kwashiorkor} Severe malnutrition in which too little protein and other nutrients cause swelling, poor growth, muscle loss, skin or hair changes, and weaker immunity. It can occur in children or adults.

\medterm{Kyasanur Forest Disease} An acute tick-borne viral hemorrhagic fever caused by Kyasanur Forest disease virus, a flavivirus transmitted to humans and non-human primates primarily by Haemaphysalis ticks (Haemaphysalis spinigera) in the Western Ghats region of South Asia. Following an incubation period of 3 to 8 days, it presents with sudden onset of high fever, frontal headache, conjunctival injection, severe prostration, myalgia, and hemorrhagic manifestations (such as epistaxis, hematemesis, or melena), frequently followed by a secondary biphasic stage marked by meningoencephalitis and tremors.

\synonyms
KFD; Monkey Fever; Kyasanur Forest Disease Virus Infection

\medterm{Kynurenine} A substance made when the body breaks down the amino acid tryptophan. It is involved in chemical signaling in the brain and immune system and may change during infection, inflammation, or some diseases.

\medterm{Kyphoplasty} A minimally invasive procedure that uses a balloon to create space in a collapsed vertebral body and then injects bone cement to stabilize a painful compression fracture.

\medterm{Kyphoscoliosis} A combination of sideways curvature and excessive forward rounding of the spine. Severe deformity can reduce chest movement and lung capacity, causing breathlessness or, over time, breathing failure and high blood pressure in the lung arteries.

\medterm{Kyphosis} An exaggerated forward curve of the upper spine, sometimes producing a rounded back. Causes include posture, osteoporosis and spinal fractures, growth disorders, arthritis, or inflammation. It may cause pain, weakness, or worsening curvature.

\synonyms
Round Back

\medterm{Kyrle Disease} A perforating dermatosis characterized by large, intensely itchy, keratotic papules that expel abnormal dermal material through the epidermis. It is strongly associated with diabetes and chronic kidney disease and may improve when associated illness is controlled.

\medterm{Kytococcus} A genus of gram-positive cocci found on skin and in the environment that can rarely cause opportunistic infection.

\medterm{Kytococcus Sedentarius} A skin-associated gram-positive coccus that produces enzymes linked to pitted keratolysis and can rarely cause invasive infection.
